\documentclass[12pt]{article}
\begin{document}


\section{Class Questions}
Turn each into If P then Q.
\begin{enumerate}
\item A matrix is invertible provided that its determinant is not zero.
\item For a function to be continuous, it is sufficient that it is differentiable.
\item For a function to be integrable, it is necessary that it is continuous.
\item A function is rational if it is a polynomial.
\item An integer is divisible by 8 only if it is divisible by 4.
\item Whenever a surface has only one side, it is non-orientable.
\item A series converges whenever it converges absolutely.
\item A geometric series with ratio r converges if $|r|<1$
\item A function is integrable provided the function is continuous.
\item The discriminant is negative only if the quadratic equation has no real solutions.
\item You fail only if you stop writing. (Ray Bradbury)
\item People will generally accept facts as truth only if the facts agree with what
they already believe. (Andy Rooney)
\item Whenever people agree with me I feel I must be wrong. (Oscar Wilde)
\end{enumerate}
\subsection{Breakout Sessions}

If $m\in \mathbf{Z}$ is even, then $m^2$ is even.

Proof? We know that by the definition of an even integer, there exists some $n\in\mathbf{Z}$ such that $m=2n$
So if $m^2$ is an even integer, then $m^2=2n$, if so, $m=\sqrt{2n}$, but $\sqrt{2}$ is never an integer, so our supposition is false.  

Where did we go wrong?



\pagebreak
\section{Proof Examples}

Proof by cases:
Prove Theorem: If q is not divisible by 3, then $q^2$ divided by 3 leaves 1 remaining.

q is not divisible by 3, so it either leaves 1 or 2 remaining.

case 1:
$q=3k+1$ for some $k\in\mathbf{Z}$
so
$$q^2=(3k+1)^2=9k^2+6k+1$$
$$=3(k^2+2k)+1$$
So $q^2$ is divisible by 3 with a remainder of 1

case 2:
$q=3k+2$ for some $k\in\mathbf{Z}$
so
$$q^2=(3k+2)^2=9k^2+12k+4$$
$$=9k^2+12k+3+1$$
$$=3(3k^2+4k+1)+1$$

So $q^2$ is divisible by 3 with a remainder of 1

So in either case $q^2$ is divisible by 3 with a remainder of 1
\subsection{Proof Examples}
Prove: If $q^2$ is divisible by 3, so is q.

Use contrapositive, if q is not divisible by 3 then $q^2$ is not divisible by 3.
Theorem above tells us that.


\end{document}

